% PURPOSE: Define phases as ETS.K_EA.v1.1 predicates and the induced partition of Omega(K_EA).
% TARGET: referees and formal-methods readers.
% LANGUAGE MODE: strictly formal.
% EVIDENCE PROFILE: none (definitions only; predicates are direct ETS transcription).
% ROLE: definitional firewall: phase identification only, no intervention semantics.

\section{Phase definitions as \texorpdfstring{$\Theta$}{Theta}-induced partitions}
\label{sec:phases}

\begin{definition}[Phase predicate (ETS-faithful)]
Let $s \in \Omega(K_{EA})$ be an admissible state. A phase predicate is a Boolean-valued
function
\[
\pi : \Omega(K_{EA}) \rightarrow \{\text{true}, \text{false}\},
\]
defined exclusively in terms of variables, components, and predicates already
present in \artifact{ETS.K_EA.v1.1}.
\end{definition}

\begin{definition}[Phase logic and precedence (ETS)]
The ETS phase logic is the totally ordered set
\[
\texttt{STOP} \succ \texttt{COLLAPSE} \succ \texttt{INERTIA} \succ \texttt{SUCCESS},
\]
where $\succ$ denotes strict precedence. If multiple phase predicates evaluate
to true for a given state $s$, the phase with highest precedence is assigned.
\end{definition}

\begin{definition}[ETS phase predicates]
The phase predicates are transcribed directly from
\artifact{ETS.K_EA.v1.1} and evaluated pointwise for $s \in \Omega(K_{EA})$:
\begin{align}
\pi_{\texttt{STOP}}(s)
&:= \bigl(f_{\text{exist}}(s) > 0\bigr)
    \,\lor\, \bigl(f_{\text{stop}}(s) > 0\bigr)
    \,\lor\, \artifact{observability\_lost}(s),
    \label{eq:pi-stop} \\[4pt]
\pi_{\texttt{COLLAPSE}}(s)
&:= \bigl(f_{\text{collapse}}(s) > 0\bigr)
    \,\land\, \neg \pi_{\texttt{STOP}}(s),
    \label{eq:pi-collapse} \\[4pt]
\pi_{\texttt{INERTIA}}(s)
&:= \bigl(f_{\text{inertia}}(s) > 0\bigr)
    \,\land\, \neg\!\bigl(\pi_{\texttt{STOP}}(s)
                          \lor \pi_{\texttt{COLLAPSE}}(s)\bigr),
    \label{eq:pi-inertia} \\[4pt]
\pi_{\texttt{SUCCESS}}(s)
&:= \bigwedge_{k} \bigl(f_k(s) \le 0\bigr),
    \label{eq:pi-success}
\end{align}
where
\[
f_k \in \{f_{\text{exist}}, f_{\text{stab}}, f_{\text{inertia}},
          f_{\text{collapse}}, f_{\text{stop}}\}
\]
are the ETS-defined structural-tension components.
\end{definition}

\begin{definition}[Phase classes as subsets of $\Omega(K_{EA})$]
Each phase predicate induces a subset of the admissible state region:
\[
\Omega_{\texttt{PHASE}}
:= \{\, s \in \Omega(K_{EA}) \mid \pi_{\texttt{PHASE}}(s)=\text{true} \,\},
\quad
\texttt{PHASE} \in
\{\texttt{SUCCESS}, \texttt{INERTIA}, \texttt{COLLAPSE}, \texttt{STOP}\}.
\]
The effective classification map
\[
\Pi : \Omega(K_{EA}) \rightarrow
\{\texttt{SUCCESS}, \texttt{INERTIA}, \texttt{COLLAPSE}, \texttt{STOP}\}
\]
assigns to each state $s$ the unique highest-precedence phase whose predicate
is satisfied at $s$.
\end{definition}

\begin{remark}[Diagnostic meaning only]
Phase membership $\Pi(s)$ is a diagnostic classification of a state under
ETS predicates. No intervention objective, control policy, recovery claim,
or optimization target is inferred from $\Pi$ within this paper.
\end{remark}

\begin{definition}[Allowed phase transitions (ETS)]
ETS specifies a phase-level allowance relation $\Rightarrow$ given by
\begin{align}
\texttt{SUCCESS}  &\Rightarrow \{\texttt{INERTIA}, \texttt{COLLAPSE}, \texttt{STOP}\}, \\
\texttt{INERTIA}  &\Rightarrow \{\artifact{SUCCESS\_via\_external\_Psi},
                              \texttt{COLLAPSE}, \texttt{STOP}\}, \\
\texttt{COLLAPSE} &\Rightarrow \{\texttt{STOP}\}, \\
\texttt{STOP}     &\Rightarrow \varnothing .
\end{align}
This relation specifies admissibility only. In particular, the label
\artifact{SUCCESS\_via\_external\_Psi} is treated as an ETS-defined symbol and
does not encode an actionable mechanism or prescription.
\end{definition}
